\begin{figure}[H]
\centering
\begin{tikzpicture}[>=Latex,scale=0.98]
    % Tren y puntos relevantes al inicio de T_1
    \draw[black,very thick,rounded corners=2pt,fill=gray!8]
        (-4.8,-0.1) rectangle (4.8,1.15);
    \draw[->,blue!75!black,very thick]
        (1.2,2.48) -- (3.8,2.48)
        node[midway,above] {$\vec v$};

    \coordinate (M0) at (-2.15,0.45);
    \coordinate (F0) at (4.2,0.78);
    \coordinate (E) at (1.25,0.62);

    \draw[gray!70!black,thick] (M0) -- ++(0,0.55);
    \node[font=\small,align=center] at (-2.15,-0.72)
        {posición inicial\\del vagón $m$};
    \fill[red!75!black] (F0) circle (0.1);
    \node[red!75!black,font=\small,align=center] at (4.05,-0.72)
        {reflexión en\\el frente};

    % Desplazamiento del vagón m y recorrido de la luz
    \draw[->,blue!75!black,very thick]
        (M0) -- (E)
        node[midway,below=5pt] {$vT_1$};
    \draw[->,red!75!black,very thick]
        (F0) -- (E)
        node[midway,above=5pt] {$cT_1$};
    \draw[black,thick] (E) circle (0.2);
    \fill[black] (E) circle (0.07);
    \node[font=\scriptsize,align=center,fill=white,inner sep=1pt]
        at (1.25,0.16) {encuentro en $m$};

    % Separación geométrica fL
    \draw[gray!65,densely dotted] (-2.15,1.0) -- (-2.15,2.02);
    \draw[gray!65,densely dotted] (4.2,0.9) -- (4.2,2.02);
    \draw[<->,black,thick]
        (-2.15,1.88) -- (4.2,1.88)
        node[midway,fill=white,inner sep=2pt]
        {$fL=cT_1+vT_1$};
\end{tikzpicture}
\caption{Geometría del encuentro durante $T_1$. La luz recorre $cT_1$ hacia atrás mientras el vagón $m$ avanza $vT_1$; ambos recorridos suman la separación inicial $fL$.}
\label{fig:tren-luz-t1}
\end{figure}